\documentclass{article}
\usepackage[utf8]{inputenc}
\usepackage{amsmath}
\usepackage{booktabs}
\usepackage{multirow}
\usepackage{graphicx}
\usepackage{geometry}
\geometry{a4paper, margin=1in}

\begin{document}

\section*{Wing Planform and Twist Trade Study}

\subsection*{1. Methodology \& Formulation}
To determine the optimal wing configuration using the \textbf{S1223} airfoil—evaluating various taper ratios ($0.6$, $0.8$, and $1.0$) -a Weighted Decision Matrix was implemented at a velocity of $V = 20.0\,\text{m/s}$ using VLM2. Each aerodynamic parameter $i$ is assigned a weight $W_i$ based on mission priorities such that:
\begin{equation}
    \sum_{i=1}^{n} W_i = 1.00
\end{equation}

Each configuration $j$ receives a discrete score $s_{ij} \in \{-1, 0, 1\}$ representing worst, average/baseline, and best performance, respectively. The total score $S_j$ for configuration $j$ is calculated as:
\begin{equation}
    S_j = \sum_{i=1}^{n} \left( W_i \cdot s_{ij} \right)
\end{equation}

\subsection*{2. Extracted Aerodynamic Parameters}
The comparative aerodynamic data extracted from XFLR5 simulations for the evaluated wing variations are summarized in Table~\ref{tab:extracted_data}.

\begin{table}[h!]
\centering
\caption{Extracted Wing Aerodynamic and Stability Data ($V = 20.0\,\text{m/s}$)}
\label{tab:extracted_data}
\begin{tabular}{lccccc}
\toprule
\textbf{Configuration} & \textbf{Max $C_L/C_D$} & \textbf{Max $\sqrt{C_L^3}/C_D^2$} & \textbf{$C_{m}$ Slope ($\partial C_m / \partial \alpha$)} & \textbf{Max $C_L$} \\
\midrule
S1223W-1.0T (Untwisted)   & 11.45 & 11.62 & Moderate Negative & 1.82 \\
S1223W-0.8T (Untwisted)   & 12.80 & 13.01 & Good Negative     & 1.85 \\
S1223W-0.6T (Untwisted)   & 13.90 & 14.15 & Strong Negative   & 1.88 \\
\textbf{S1223W-0.6T-7Twisted} & \textbf{14.25} & \textbf{14.48} & \textbf{Optimal Negative} & \textbf{1.88} \\
\bottomrule
\end{tabular}
\end{table}

\subsection*{3. Aerodynamic Performance Comparison (Visuals)}
Figure~\ref{fig:wing_graphs} illustrates the comparative XFLR5 graphs for the different taper ratios and twist configurations, showing the distinct improvements in lift-to-drag ratios, induced drag characteristics, and pitching moment stability slopes.

\begin{figure}[h!]
\centering
% قم بتعديل مسار وصيغة اسم الصورة هنا بما يتناسب مع ملفك
\includegraphics[width=0.9\textwidth]{wing graphs.png}
\caption{XFLR5 Aerodynamic Comparison Graphs for Wing Variations ($0.6T$, $0.8T$, $1.0T$)}
\label{fig:wing_graphs}
\end{figure}

\subsection*{4. Weight Distribution \& Evaluation Criteria}
\begin{itemize}
    \item \textbf{Peak Aerodynamic Efficiency ($(C_L/C_D)_{\max}$, $W_1 = 0.30$):} Crucial for maximizing glide range and overall cruise performance. Higher is better.
    \item \textbf{Climb/Endurance Metric ($\max \sqrt{C_L^3}/C_D^2$, $W_2 = 0.25$):} Dictates power required and sustained climb capabilities. Higher is better.
    \item \textbf{Longitudinal Stability ($C_m$ Slope \& Trim, $W_3 = 0.25$):} Ensures static longitudinal stability and smooth pitching behavior. Stronger negative slope is better.
    \item \textbf{Lift Capability ($C_{L,\max}$, $W_4 = 0.20$):} Influences stall margins and slow-speed control. Higher is better.
\end{itemize}

\subsection*{5. Trade Study Decision Matrix}
Table~\ref{tab:trade_matrix} presents the raw ratings, assigned weights, and weighted score calculations across the core candidate configurations.

\begin{table}[h!]
\centering
\caption{Wing Configuration Trade Study Matrix}
\label{tab:trade_matrix}
\begin{tabular}{lccccccr}
\toprule
\multirow{2}{*}{\textbf{Parameter}} & \multirow{2}{*}{\textbf{Weight ($W_i$)}} & \multicolumn{3}{c}{\textbf{Scores ($s_{ij}$)}} & \multicolumn{3}{c}{\textbf{Weighted Scores ($W_i \cdot s_{ij}$)}} \\
\cmidrule(lr){3-5} \cmidrule(lr){6-8}
& & \textbf{1.0T} & \textbf{0.8T} & \textbf{0.6T-7Twisted} & \textbf{1.0T} & \textbf{0.8T} & \textbf{0.6T-7Twisted} \\
\midrule
$(L/D)_{\max}$          & 0.30 & -1 & 0  & +1 & -0.30 & 0.00  & 0.30 \\
$\sqrt{C_L^3}/C_D^2$   & 0.25 & -1 & 0  & +1 & -0.25 & 0.00  & 0.25 \\
Stability ($C_m$ slope) & 0.25 & -1 & 0  & +1 & -0.25 & 0.00  & 0.25 \\
$C_{L,\max}$            & 0.20 & -1 & 0  & +1 & -0.20 & 0.00  & 0.20 \\
\midrule
\textbf{Total Score ($S_j$)} & \textbf{1.00} & -- & -- & -- & \textbf{-1.00} & \textbf{0.00} & \textbf{+1.00} \\
\bottomrule
\end{tabular}
\end{table}

\subsection*{6. Detailed Mathematical Calculations}
The final weighted total scores are computed explicitly to justify the engineering selection:

\begin{align}
    S_{\text{1.0T}}               &= (0.30 \times -1) + (0.25 \times -1) + (0.25 \times -1) + (0.20 \times -1) = \mathbf{-1.00} \\
    S_{\text{0.8T}}               &= (0.30 \times 0)  + (0.25 \times 0)  + (0.25 \times 0)  + (0.20 \times 0)  = \mathbf{0.00}  \\
    S_{\text{0.6T-7Twisted}}      &= (0.30 \times 1)  + (0.25 \times 1)  + (0.25 \times 1)  + (0.20 \times 1)  = \mathbf{+1.00}
\end{align}

\subsection*{7. Selection Conclusion}
According to the trade study results and visual verification from the aerodynamic graphs, the tapered wing configuration with a $0.6$ taper ratio combined with a $7^\circ$ washout twist (\textbf{S1223W-0.6T-7Twisted}) achieves the highest total score ($S = +1.00$). This configuration is conclusively recommended as the optimal design because it successfully combines superior aerodynamic efficiency ($L/D$), optimized induced drag characteristics, excellent static longitudinal stability, and a safer stall progression.

\end{document}